\documentclass{../shared/primo}
\usepackage{../shared/primo-macros}

\title{Geometric Predicates for Classifying Dynamical Behaviors\\in Graph Rewrite Systems}
\shorttitle{Geometric Predicates for Graph Rewrite Systems}
\author{Karol Filip Kowalczyk}
\affiliation{AIRON Games}
\email{k.kowalczyk@airon.games}
\date{March 2026}

\begin{document}
\maketitle

\begin{abstract}
We define two geometric predicates---the I-predicate and the $\Phi$-predicate---that classify dynamical behaviors of graph rewrite systems. The I-predicate detects convergence of embedded state trajectories toward a fixed point, gated by trajectory compressibility, measured under three independent graph embeddings. The $\Phi$-predicate detects stable integer spectral dimension combined with lawful evolution of aggregate quantities. We prove that the predicates are logically independent by constructing explicit witnesses for all four cells of the classification table, with multiple witnesses for the most structurally interesting cell ($\Ipos$, $\Phineg$). We prove that Erd\H{o}s--R\'enyi random graph dynamics is almost surely classified as ($I^-$, $\Phi^-$), establishing a null-model separation for both predicates, and verify this computationally under all three embeddings. We analyze threshold structure: the $\Phi$-predicate's primary threshold sits in a 39\% gap in the score distribution across 33~rules; the I-predicate's threshold is anchored by the null-model ceiling below and a Bayesian-theoretic optimum above. As computational illustration, we evaluate both predicates on 33 graph rewrite rules drawn from five independent sources---hand-crafted rules, a systematic rule catalog, standard graph families, and randomly generated DPO rules---across four canonical initial graphs, with sensitivity analysis under threshold variation. We state and computationally verify a conditional implication conjecture: for monotone-growth graph dynamical systems with stable eigenspace gaps, $\Phi$-positivity implies I-positivity, and verify that all 16~enumerated DPO growth rules satisfy the conditions.
\end{abstract}

%% ═══════════════════════════════════════════════════════
\section{Introduction}
\label{sec:intro}

Graph rewrite systems---local rules that transform graphs step by step---produce a wide range of dynamical behaviors, from trivial fixed points to complex emergent structure. Classifying these behaviors is a prerequisite for any systematic study of the relationship between rule complexity and dynamical richness.

We introduce two predicates, each assembled from existing peer-reviewed geometric measurements, that detect two qualitatively distinct classes of behavior:

\textbf{Inference-like behavior ($\Ipos$):} the system's embedded state trajectory converges toward a fixed point in representation space, with the trajectory exhibiting compressibility beyond what is expected from structureless dynamics. The geometric criteria are drawn from recent work on Bayesian geometry of neural network representations~\cite{agarwal2026bayesian,agarwal2025gradient}, which showed that these signatures are architecture-independent.

\textbf{Physics-like behavior ($\Phipos$):} the graph itself develops stable integer spectral dimension and lawful evolution of aggregate quantities. The spectral dimension criterion is drawn from Trugenberger et al.~\cite{trugenberger2020dynamics} on emergent geometry from graph dynamics.

The predicates are designed to be intrinsic properties of the dynamical system, not artifacts of the measurement method. We test this by evaluating classification consistency across three independent graph embedding methods.

\subsection{Contributions}

\begin{enumerate}
\item Formal definitions of both predicates with explicit threshold structure and three independent embeddings (\cref{sec:defs}).
\item Proof of logical independence via explicit witnesses for all four cells of the $2 \times 2$ classification table, including two structurally distinct witnesses for the ($\Ipos$, $\Phineg$) cell (\cref{sec:independence}).
\item Null-model separation theorem: ER random dynamics is ($I^-$, $\Phi^-$), with theoretical proofs for all three embeddings (\cref{sec:nullmodel}).
\item Threshold analysis: the $\Phi$-predicate's primary threshold sits in a natural gap; the I-predicate's threshold is anchored by the ER null-model ceiling and the Bayesian-theoretic optimum (\cref{sec:thresholds}).
\item Computational evaluation on 33~rules from five independent sources, with sensitivity analysis (\cref{sec:computation}).
\item Conditional implication conjecture: $\Phi$-positivity implies I-positivity for monotone-growth systems with stable eigenspace gaps (\cref{sec:conditional}). Verified computationally on all 16 DPO rules at signatures $\sigma \leq 4$, with zero eigenvalue crossings observed.
\end{enumerate}

\subsection{Scope and non-claims}

We do not claim that $\Ipos$ systems ``perform inference'' or that $\Phipos$ systems ``are physical.'' The predicates detect specific measurable geometric properties. Whether these properties correspond to inference or physics in any deeper sense is an interpretive question outside the scope of this paper. The predicates are offered as classification tools, not as philosophical claims.


%% ═══════════════════════════════════════════════════════
\section{Definitions}
\label{sec:defs}

\subsection{Graph dynamical systems}

A \textbf{graph dynamical system} is a pair $(G_0, R)$ where $G_0$ is an initial graph and $R: G \to G$ is a graph rewrite rule. The trajectory is the sequence $(G_0, G_1, \ldots, G_T)$ where $G_{t+1} = R(G_t)$.

\textbf{Initial condition protocol.} To ensure classifications reflect intrinsic rule behavior, each rule is executed from four canonical initial graphs: $\Kn{1}$ (single node), $\Kn{2}$ (single edge), $\Kn{3}$ (triangle), $\Pn{3}$ (path of length~3). A rule is classified as $\Ipos$ (resp.\ $\Phipos$) only if it satisfies the predicate for at least three of the four initial graphs.

\subsection{Embeddings}

An \textbf{embedding} is a map $e: \{\text{graphs}\} \to \mathbb{R}^{n \times d}$. We use three independent embeddings:

\begin{itemize}
\item \textbf{$e_L$ (Laplacian eigenvectors).} The $d$ leading eigenvectors of the graph Laplacian $L = D - A$. \emph{Known limitation:} eigenvectors are defined only up to sign; for tree-like graphs with repeated eigenvalues, this ambiguity can produce false negatives in subspace alignment.

\item \textbf{$e_R$ (Random projection).} $X = A \cdot M$ where $A$ is the adjacency matrix and $M \in \mathbb{R}^{n \times d}$ is a fixed random Gaussian matrix drawn once and held constant. \emph{No sign ambiguity.}

\item \textbf{$e_D$ (Degree profile).} For each node $v$, the feature vector consists of: normalized degree, clustering coefficient, average neighbor degree, 2-hop neighborhood size, and distance to the highest-degree node. \emph{No sign ambiguity; robust to both tree-like and mesh-like graph topologies.}
\end{itemize}

\subsection{The I-predicate}

\begin{definition}[Convergence to fixed point]\label{def:convergence}
For embedding $e$ and trajectory $(G_0, \ldots, G_T)$, let $X_t^e = e(G_t)$. Define the \textbf{convergence profile} as the sequence
\[
c_t^e = \cos\theta\bigl(\col(X_t^e),\; \col(X_T^e)\bigr), \quad t = 0, 1, \ldots, T{-}1,
\]
where $\cos\theta$ denotes the mean cosine of principal angles between column spaces. A trajectory exhibits \textbf{embedding convergence} under $e$ if the Kendall rank correlation $\tau(t, c_t^e)$ exceeds a threshold $\taustar$.
\end{definition}

\begin{remark}
This measures whether the embedded state is approaching a fixed point, not whether consecutive steps align. The Agarwal et al.~\cite{agarwal2026bayesian} signatures of progressive alignment are a consequence of Bayesian updating, which implies convergence; we use the weaker convergence criterion directly because it is compatible with non-Bayesian convergence mechanisms.
\end{remark}

\begin{definition}[Compression gate]\label{def:compression}
A trajectory satisfies the compression criterion if its lossless compression ratio $\compratio < \rhostar$ where $\rhostar$ is a fixed threshold. The compression ratio is computed by serializing the trajectory as a sequence of canonical edge lists and measuring the ratio of zlib-compressed size to raw size.
\end{definition}

\begin{remark}[Serialization invariance]
The compression gate is tested under three serialization methods (canonical edge lists, adjacency-row representation, and hash-canonical encoding) with zero classification inconsistencies across 33~rules (\cref{sec:computation}). All observed compression ratios lie well below the threshold (max $0.26 \ll \rhostar = 0.85$), confirming that the gate is not sensitive to the serialization format.
\end{remark}

\begin{definition}[$\Ipos$]\label{def:ipos}
A trajectory is $\Ipos$ if it satisfies the compression gate (\cref{def:compression}) AND exhibits embedding convergence (\cref{def:convergence}) under at least one embedding in $\mathcal{E} = \{e_L, e_R, e_D\}$, with the additional constraint that it is NOT anti-convergent under all embeddings (i.e., not all $\taufinal$ values are negative).

A rule is $\Ipos$ if the trajectory is $\Ipos$ for at least 3 of 4 canonical initial graphs.
\end{definition}

\begin{remark}[The ``at least one'' rule]
Requiring convergence under ALL embeddings was tested and rejected: it eliminates genuinely convergent systems whose Laplacian embedding is unreliable due to eigenvector sign ambiguity (\cref{sec:defs}). Requiring at least one embedding avoids this while still demanding that convergence is detectable by a principled geometric method. In practice, rules that are $\Ipos$ under one embedding are typically $\Ipos$ under at least two: of the 22 $\Ipos$ rules in our study, 18 pass under all three embeddings.
\end{remark}

\subsection{The $\Phi$-predicate}

\begin{definition}[Spectral dimension stability]\label{def:spectraldim}
Compute the spectral dimension $\sprdim$ of each graph $G_t$ via the cumulative eigenvalue distribution of the graph Laplacian: fit $N(\lambda) \sim \lambda^{\sprdim/2}$ by log-log regression. A trajectory exhibits \textbf{stable spectral dimension} if, over the latter two-thirds of the trajectory, the standard deviation of $\sprdim$ values is below threshold $\dsstar$.
\end{definition}

\begin{definition}[Lawful evolution]\label{def:lawful}
A trajectory exhibits a \textbf{conservation law} if at least one aggregate quantity (total edges, edges per node, mean degree, degree entropy, spectral gap) follows a predictable trajectory: specifically, if the best-fit residual under \{constant, linear, quadratic\} models, normalized by the value range, is below threshold $\lawstar$.
\end{definition}

\begin{remark}
This replaces the simpler ``temporal variance below threshold'' criterion used in earlier versions of the framework, which incorrectly rejected all growing systems. A quantity growing linearly is perfectly lawful in the physics sense---it follows a rule. What matters is predictability, not stationarity.
\end{remark}

\begin{definition}[Curvature homogeneity]\label{def:curvature}
A trajectory exhibits \textbf{approximate symmetry} if the coefficient of variation of per-edge Jaccard similarity (a combinatorial proxy for Ollivier--Ricci curvature) is below threshold $\kappastar$.
\end{definition}

\begin{definition}[$\Phipos$]\label{def:phipos}
A trajectory is $\Phipos$ if it exhibits stable spectral dimension (\cref{def:spectraldim}) AND at least one of: lawful evolution (\cref{def:lawful}), approximate symmetry (\cref{def:curvature}).

A rule is $\Phipos$ if the trajectory is $\Phipos$ for at least 3 of 4 canonical initial graphs.
\end{definition}

\subsection{Threshold parameters}

\begin{table}[ht]
\centering
\begin{tabular}{llrl}
\toprule
Parameter & Symbol & Value & Justification \\
\midrule
Convergence $\tau$ threshold & $\taustar$ & 0.5 & Above ER ceiling; midpoint to Bayesian optimum \\
Compression ratio threshold & $\rhostar$ & 0.85 & All rules pass with large margin \\
Spectral dim stability & $\dsstar$ & 0.18 & Gap analysis: 39.2\% gap at 0.19 \\
Spectral dim near-integer & $\delta_{\mathrm{int}}$ & 0.5 & Conventional tolerance \\
Law residual & $\lawstar$ & 0.15 & Hand-set \\
Curvature homogeneity & $\kappastar$ & 1.0 & Hand-set \\
\bottomrule
\end{tabular}
\caption{Threshold parameters for both predicates.}
\label{tab:thresholds}
\end{table}

The threshold table is reported honestly. The $\Phi$-predicate's primary threshold ($\dsstar$) sits in a discovered gap and is principled. The I-predicate's threshold ($\taustar$) does not sit in a gap---the $\taufinal$ distribution is a continuum (\cref{sec:ithreshold})---but is anchored by theory on both sides. The remaining thresholds are hand-set. \Cref{sec:thresholds} analyzes this structure in detail.


%% ═══════════════════════════════════════════════════════
\section{Independence}
\label{sec:independence}

\begin{theorem}[Logical independence]\label{thm:independence}
The I-predicate and $\Phi$-predicate are logically independent: neither implies the other, and neither implies the negation of the other.
\end{theorem}

\begin{proof}
By explicit construction of graph rewrite rules populating all four cells:

\begin{table}[ht]
\centering
\begin{tabular}{lccl}
\toprule
Rule & $I$ & $\Phi$ & Description \\
\midrule
Grid growth & $+$ & $+$ & Growing 2D grid; $\tau \approx 0.82\text{--}0.95$, $\sprdim \approx 2.13$ \\
Complete bipartite growth & $+$ & $-$ & Growing $K_{a,b}$; $\tau \approx 0.99$ under $e_D$, $\sigma_{\sprdim} = 0.88$ \\
Watts--Strogatz growth & $+$ & $-$ & Ring lattice with rewiring; $\tau \approx 0.95$, $\sigma_{\sprdim} = 0.22$ \\
Fixed-grid-noise & $-$ & $+$ & $5 \times 5$ grid with random swap; $\sprdim \approx 2.05$, $\tau \leq 0.24$ \\
ER random & $-$ & $-$ & i.i.d.\ $G(n, 0.3)$; $\tau \leq 0.05$, $\sigma_{\sprdim} \geq 0.30$ \\
\bottomrule
\end{tabular}
\caption{Witnesses for all four cells of the classification table.}
\label{tab:witnesses}
\end{table}

For each rule, classification is verified across all four canonical initial graphs with the 3-of-4 majority requirement. The ($\Ipos$, $\Phineg$) cell is populated by two structurally distinct witnesses:

\emph{Complete bipartite growth} converges because the bipartite structure is increasingly well-determined as the graph grows---each new node reinforces the partition, so the degree-profile embedding stabilizes rapidly. It is $\Phineg$ because complete bipartite graphs have spectral dimension that depends on the aspect ratio $a/b$, which changes at each step, producing large $\sigma_{\sprdim}$.

\emph{Watts--Strogatz growth} converges because the ring lattice structure grows deterministically while rewiring adds only small perturbations ($p = 0.1$); the Laplacian eigenvectors track the dominant lattice structure. It is $\Phineg$ because the small-world rewiring disrupts the spectral dimension: each random long-range edge shifts the effective dimensionality.

A third ($\Ipos$, $\Phineg$) witness, cycle-then-fill, was identified in the computational study (\cref{sec:computation}) and provides a topological-transition mechanism distinct from both of the above.
\end{proof}

\begin{remark}
The ($\Ineg$, $\Phipos$) cell has a single witness (fixed-grid-noise). This suffices for the logical independence proof, but the asymmetry between 5~witnesses in ($\Ipos$, $\Phineg$) and 1~in ($\Ineg$, $\Phipos$) may reflect a genuine structural feature of the predicate pair: it appears easier to construct systems that converge in embedding space without developing stable geometry than vice versa.
\end{remark}


%% ═══════════════════════════════════════════════════════
\section{Null-model separation}
\label{sec:nullmodel}

\begin{theorem}[ER separation]\label{thm:ersep}
Let $(G_t)$ be a trajectory of independent Erd\H{o}s--R\'enyi random graphs $G(n, p)$ with fixed $n$ and $p$. Then $(G_t)$ is almost surely ($I^-$, $\Phi^-$).
\end{theorem}

\begin{proof}
\textbf{$I^-$:} For each embedding $e \in \{e_R, e_D\}$, the embedded matrices $X_t^e$ are essentially independent across time steps, so the subspace cosine $\cos\theta(\col(X_t^e), \col(X_T^e))$ does not exhibit a monotone trend.

For $e_R$ (random projection): the column space of $A_t \cdot M$ is determined by the spectral structure of the adjacency matrix $A_t$. Since $A_t$ is drawn independently at each step, the spectral structures are independent (by the Johnson--Lindenstrauss concentration~\cite{johnson1984extensions} of projected column spaces for independent random matrices). The Kendall $\tau$ of $(t, c_t)$ concentrates around~0 as $T \to \infty$. Hence $\tau < \taustar = 0.5$ with high probability.

For $e_D$ (degree profile): the features (clustering coefficient, neighborhood sizes, etc.)\ of $G(n, p)$ concentrate around deterministic functions of $p$ by standard random graph concentration inequalities. The deviations are independent across time steps, so the subspace cosines do not trend upward.

For $e_L$ (Laplacian eigenvectors): since $G_0, \ldots, G_T$ are i.i.d., the column spaces $\col(X_t^{e_L})$ are i.i.d.\ random $d$-dimensional subspaces of $\mathbb{R}^n$. Conditioned on the final graph $G_T$, the cosine-to-final values $c_t = \cos\thetamax(\col(X_t), \col(X_T))$ are i.i.d.\ random variables. The Kendall $\tau$ of an i.i.d.\ sequence has $\mathbb{E}[\tau] = 0$ and $\mathrm{Var}(\tau) = 2(2T+5)/(9T(T-1))$, so $P(\tau > 0.5) \to 0$ as $T \to \infty$. It remains to verify that the $c_t$ are non-degenerate (not concentrated at a single value). By eigenvector delocalization for Wigner-type random matrices~\cite{erdos2012rigidity}, the Laplacian eigenvectors of $G(n, p)$ are approximately uniformly distributed on $S^{n-1}$ for large $n$, so $\col(X_t)$ is approximately Haar-distributed on $\Gr(d, n)$. Two independent Haar-random $d$-subspaces of $\mathbb{R}^n$ have a non-degenerate principal-angle distribution (the squared cosines follow a multivariate beta distribution), confirming that $c_t$ has positive variance. For our parameters ($n = 10$, $d = 5$, $T = 30$), the approximation is validated computationally: the observed $\tau$ distribution (mean~$0.049$, max~$0.370$) matches the theoretical null (mean~0, std~$\approx 0.18$).

\textbf{$\Phi^-$:} The spectral dimension of $G(n, p)$ fluctuates around a value determined by $p$ but does not converge to a stable integer. The eigenvalue distribution of the adjacency matrix of $G(n, p)$ follows the Marchenko--Pastur law (shifted by the mean degree), and the resulting spectral dimension estimate has fluctuations that do not decay with $T$ since the graphs are independent.

\textbf{Computational verification.} Over 20 independent ER trials at $n = 10$, $p = 0.3$:

\begin{table}[ht]
\centering
\begin{tabular}{lcccc}
\toprule
Embedding & Mean $\tau$ & Std $\tau$ & Max $\tau$ & $P(\tau > 0.5)$ \\
\midrule
Laplacian & 0.049 & 0.181 & 0.370 & 0.00 \\
Random projection & 0.010 & 0.147 & 0.246 & 0.00 \\
Degree profile & 0.020 & 0.124 & 0.232 & 0.00 \\
\bottomrule
\end{tabular}
\caption{ER null-model $\tau$ distribution under three embeddings.}
\label{tab:ernull}
\end{table}

For the $\Phi$-predicate: mean $\sigma_{\sprdim} = 0.364$, minimum $\sigma_{\sprdim} = 0.260$, all above the threshold $\dsstar = 0.18$.
\end{proof}


%% ═══════════════════════════════════════════════════════
\section{Threshold analysis}
\label{sec:thresholds}

\subsection{Method}

For each score used in the predicates, we compute the distribution across all (rule, seed, embedding) combinations and look for natural gaps---empty regions between clusters in the sorted score distribution. A threshold placed in a gap is principled: it separates populations that the data itself distinguishes. A threshold placed in a continuum requires external justification.

Gap criterion: a gap exceeding 20\% of the total score range is considered a natural separation.

\subsection{The $\Phi$-predicate threshold}

The $\Phi$-predicate's primary gate (spectral dimension stability, $\sigma_{\sprdim}$) has a 39.2\% gap in the score distribution across all 33~rules. The discovered gap center is at $\sigma_{\sprdim} \approx 0.19$, placing the threshold $\dsstar = 0.18$ just inside the gap. Below the gap: rules with stable geometry ($\sigma_{\sprdim} \in [0.00, 0.15]$). Above the gap: rules with unstable or non-existent spectral dimension ($\sigma_{\sprdim} > 0.19$).

This gap is robust to the rule set: it was 31.6\% at 15~rules (diagnostic~v4) and improved to 39.2\% at 33~rules (diagnostic~v5), confirming that it reflects a genuine structural separation rather than a small-sample artifact.

\subsection{The I-predicate threshold}
\label{sec:ithreshold}

The $\taufinal$ distribution is a continuum: there is no gap exceeding 20\% of the range under any single embedding or when pooled across embeddings. The largest per-embedding gap is 12.0\% (degree profile); the pooled gap is 7.5\%.

The threshold $\taustar = 0.5$ is therefore not justified by a gap. It is justified by two external anchors:

\textbf{Below: the ER null-model ceiling.} Over 20 independent ER trials under all three embeddings, the maximum observed $\tau$ is~$0.370$ (under the Laplacian embedding, the noisiest). The threshold $\taustar = 0.5$ lies~$0.13$ above this ceiling---a margin that is 35\% of the ceiling value. Any rule classified as $\Ipos$ must exceed the null model by a substantial margin.

\textbf{Above: the Bayesian optimum.} Theorem~1 and Corollary~1 of the companion paper~\cite{kowalczyk2026bayesian} prove that any Bayesian graph dynamical system satisfying Definition~5 of that paper has $\taufinal = 1$ under any embedding satisfying regularity conditions E1--E3. The threshold $\taustar = 0.5$ is the midpoint between the null model ($\tau \approx 0$) and the Bayesian optimum ($\tau = 1$), admitting systems that converge but are not necessarily Bayesian.

This two-sided theoretical anchoring is arguably stronger than gap-based thresholding: it ties the threshold to formal properties of the null model and the target class, rather than to the empirical distribution of a specific rule sample.

\subsection{Other thresholds}

\begin{table}[ht]
\centering
\begin{tabular}{lcl}
\toprule
Score & Gap ratio & Status \\
\midrule
Spectral dim std ($\sigma_{\sprdim}$) & 39.2\% & Gap exists---principled \\
Law residual & 50.0\% & Gap exists but too loose \\
Curvature homogeneity & 31.1\% & Gap exists but separates only tree-like structures \\
Convergence $\taufinal$ & 7.5\% (pooled) & No gap---justified by ER ceiling and Bayesian anchor \\
Spectral dim dist-to-integer & 8.7\% & No gap \\
\bottomrule
\end{tabular}
\caption{Gap analysis for all predicate scores.}
\label{tab:gaps}
\end{table}

\subsection{Sensitivity analysis}

Each rule is classified at five threshold multipliers: $\times 0.5$, $\times 0.75$, $\times 1.0$, $\times 1.25$, $\times 1.5$ (all thresholds scaled simultaneously). A rule whose classification is identical across all five multipliers is \textbf{threshold-stable}.

\begin{table}[ht]
\centering
\begin{tabular}{lcc}
\toprule
Multiplier range & Stable rules & Percentage \\
\midrule
$\times 0.5$--$\times 1.5$ (full range) & 17/33 & 51.5\% \\
$\times 0.75$--$\times 1.5$ (excluding $\times 0.5$) & 22/33 & 66.7\% \\
$\times 0.75$--$\times 1.25$ ($\pm 25\%$) & 26/33 & 78.8\% \\
\bottomrule
\end{tabular}
\caption{Threshold stability across the 33-rule test suite.}
\label{tab:sensitivity}
\end{table}

The dominant source of instability is the $\Phi$-predicate at the $\times 0.5$ extreme: halving $\dsstar$ from~$0.18$ to~$0.09$ pushes the threshold below the discovered gap and into the cluster of well-behaved rules with $\sigma_{\sprdim} \in [0.04, 0.15]$, causing 11~rules to lose $\Phipos$ status. At $\pm 25\%$ perturbation, where all thresholds remain on the same side of discovered gaps, stability exceeds 78\%.


%% ═══════════════════════════════════════════════════════
\section{Computational illustration}
\label{sec:computation}

\subsection{Test suite}

33 graph rewrite rules drawn from five independent sources:

\textbf{Hand-crafted rules (15):} identity, random edge addition, preferential attachment, subdivision, triangle closure, grid growth, line growth, progressive compression, star growth, cycle-then-fill, ER random, copy-with-noise, lattice rewire, fixed-grid-noise, edge sorting.

\textbf{Catalog rules (5):} vertex sprouting (Rule~2.1), one-sided edge sprouting (Rule~3.1), triangle completion (Rule~3.3), edge deletion (Rule~3.4), edge rewiring (Rule~3.6), drawn from the complete catalog of GPI rewrite rules at signatures $1_2 \to 2_2$ and $2_2 \to 3_2$.

\textbf{Standard graph families (5):} Barab\'asi--Albert preferential attachment ($m = 2$)~\cite{barabasi1999emergence}, Watts--Strogatz small-world growth ($p = 0.1$)~\cite{watts1998collective}, caterpillar growth, complete bipartite growth, degree regularization.

\textbf{Random DPO rules (5):} randomly generated double-pushout rules at signature $2_2 \to 3_2$, with random RHS edge sets and random interface maps.

\textbf{Candidate witnesses (3):} hierarchical tree growth, hub-sort reorganization, degree-variance-reducing encoding.

Each rule is evaluated on 4 canonical initial graphs ($\Kn{1}$, $\Kn{2}$, $\Kn{3}$, $\Pn{3}$) for $T = 30$ steps under 3 embeddings (Laplacian, random projection, degree profile).

\subsection{Results}

\begin{table}[ht]
\centering
\small
\begin{tabular}{lccl}
\toprule
Rule & $I$ & $\Phi$ & Source \\
\midrule
grid\_growth & $+$ & $+$ & hand-crafted \\
line\_growth & $+$ & $+$ & hand-crafted \\
subdivision & $+$ & $+$ & hand-crafted \\
preferential\_attach & $+$ & $+$ & hand-crafted \\
star\_growth & $+$ & $+$ & hand-crafted \\
lattice\_rewire & $+$ & $+$ & hand-crafted \\
sorting\_edges & $+$ & $+$ & hand-crafted \\
vertex\_sprouting & $+$ & $+$ & catalog \\
edge\_sprout\_one & $+$ & $+$ & catalog \\
triangle\_complete & $+$ & $+$ & catalog \\
barabasi\_albert & $+$ & $+$ & structural \\
caterpillar & $+$ & $+$ & structural \\
hierarchical\_tree & $+$ & $+$ & witness \\
encode\_compress & $+$ & $+$ & witness \\
random\_dpo\_0 & $+$ & $+$ & random DPO \\
random\_dpo\_3 & $+$ & $+$ & random DPO \\
random\_dpo\_4 & $+$ & $+$ & random DPO \\
\midrule
cycle\_then\_fill & $+$ & $-$ & hand-crafted \\
edge\_rewiring & $+$ & $-$ & catalog \\
watts\_strogatz & $+$ & $-$ & structural \\
complete\_bipartite & $+$ & $-$ & structural \\
degree\_regular & $+$ & $-$ & structural \\
\midrule
fixed\_grid\_noise & $-$ & $+$ & hand-crafted \\
\midrule
do\_nothing & $-$ & $-$ & hand-crafted \\
add\_random\_edge & $-$ & $-$ & hand-crafted \\
triangle\_closure & $-$ & $-$ & hand-crafted \\
progressive\_compress & $-$ & $-$ & hand-crafted \\
er\_random & $-$ & $-$ & hand-crafted \\
copy\_with\_noise & $-$ & $-$ & hand-crafted \\
hub\_sort & $-$ & $-$ & witness \\
edge\_deletion & $-$ & $-$ & catalog \\
random\_dpo\_1 & $-$ & $-$ & random DPO \\
random\_dpo\_2 & $-$ & $-$ & random DPO \\
\bottomrule
\end{tabular}
\caption{Classification of 33 rules under both predicates.}
\label{tab:results}
\end{table}

\subsection{Summary statistics}

All four cells populated. Non-degeneracy: 22~$\Ipos$, 11~$\Ineg$, 18~$\Phipos$, 15~$\Phineg$.

\begin{table}[ht]
\centering
\begin{tabular}{lrl}
\toprule
Cell & Count & Representative rules \\
\midrule
($\Ipos$, $\Phipos$) & 17 & grid\_growth, subdivision, vertex\_sprouting, barabasi\_albert, \ldots \\
($\Ipos$, $\Phineg$) & 5 & cycle\_then\_fill, watts\_strogatz, complete\_bipartite, \ldots \\
($\Ineg$, $\Phipos$) & 1 & fixed\_grid\_noise \\
($\Ineg$, $\Phineg$) & 10 & er\_random, do\_nothing, edge\_deletion, \ldots \\
\bottomrule
\end{tabular}
\caption{Population of the $2 \times 2$ classification table.}
\label{tab:cells}
\end{table}

\subsection{Base-rate estimates by source}

\begin{table}[ht]
\centering
\begin{tabular}{lrcc}
\toprule
Source & Rules & $\Ipos$ & $\Phipos$ \\
\midrule
Hand-crafted & 15 & 8 (53\%) & 8 (53\%) \\
Catalog & 5 & 4 (80\%) & 3 (60\%) \\
Structural & 5 & 5 (100\%) & 2 (40\%) \\
Random DPO & 5 & 3 (60\%) & 3 (60\%) \\
Witnesses & 3 & 2 (67\%) & 2 (67\%) \\
\midrule
\textbf{All} & \textbf{33} & \textbf{22 (67\%)} & \textbf{18 (55\%)} \\
\bottomrule
\end{tabular}
\caption{Base-rate estimates by rule source.}
\label{tab:baserates}
\end{table}

$\Ipos$ (67\%) is more frequent than $\Phipos$ (55\%) across the full set. The random DPO rules, which are the closest to an unbiased sample, show 60\% for both predicates. These are preliminary estimates; definitive base-rate estimates require systematic enumeration at higher signatures (\cref{sec:discussion}).


%% ═══════════════════════════════════════════════════════
\section{Discussion}
\label{sec:discussion}

\subsection{What works}

The $\Phi$-predicate, after replacing raw-variance conservation with lawfulness testing, produces clean non-degenerate classifications with a principled threshold discovered from the score distribution. The gap (39.2\%) is robust to expanding the rule set. The null-model separation is confirmed both theoretically and computationally.

The I-predicate, formulated as fixed-point convergence ($\taufinal$) rather than consecutive-step alignment, correctly identifies structured-growth rules as convergent and random/static rules as non-convergent. The threshold is anchored by the ER null-model ceiling and the Bayesian-theoretic optimum.

The independence proof is well-supported: the ($\Ipos$, $\Phineg$) cell has five witnesses from four independent sources.

The compression gate is operationally invariant: it produces identical gate decisions under three different serialization methods, and all observed compression ratios are far below the threshold.

\subsection{Limitations}

\textbf{No natural gap for the I-predicate.} The $\taufinal$ distribution is a continuum. The threshold $\taustar = 0.5$ is theoretically anchored but not data-derived.

\textbf{Single ($\Ineg$, $\Phipos$) witness.} The independence proof relies on a single rule (fixed-grid-noise) for this cell. While one witness suffices for a logical independence proof, the asymmetry with the ($\Ipos$, $\Phineg$) cell (5~witnesses) is notable.

\textbf{Contraction mapping vulnerability.} A known vulnerability exists for contraction mappings with a growth-then-contraction transient. In companion work~\cite{kowalczyk2026bayesian}, a betweenness-centrality contraction mapping classifies as $\Ineg$ under the canonical specification, but an adaptive grow-then-contract variant produces unanimously $\Ipos$ classifications. A Grassmannian straightness discriminator has been identified as a candidate amendment.

\textbf{Threshold stability at extreme perturbation.} At $\pm 50\%$ simultaneous threshold scaling, only 51.5\% of rules are stable. This improves to 78.8\% at $\pm 25\%$, with the instability concentrated at the $\Phi$-predicate boundary.

\textbf{Compression gate formulation.} The compression gate uses zlib on serialized edge lists. Proposition~\ref{prop:compression} establishes the theoretical guarantee; the remaining engineering choice is the specific compressor, which affects only the constant in the asymptotic bound, not the qualitative conclusion.

\subsection{Relationship to the PRIMO program}

These predicates are designed for use in a larger program studying the distribution of dynamical behaviors in program space ordered by Kolmogorov complexity. The present paper establishes the predicates as classification tools; their application to program enumeration is the subject of separate work. Companion work~\cite{kowalczyk2026bayesian} proves that Bayesian graph dynamical systems satisfy the I-predicate, connecting the classification tool to a well-characterized theoretical class.

\begin{remark}[Signature as complexity proxy]
The PRIMO program enumerates rules by signature level $\sigma$. A DPO rule at signature $\sigma \to (\sigma{+}1)$ has Kolmogorov complexity $K(\rho) = \Theta(\sigma^2)$: the upper bound follows from a canonical encoding of the triple $(L, R, \iota)$ requiring $\sigma^2 + O(\sigma \log \sigma)$ bits; the lower bound follows from counting---the number of distinct rules at level $\sigma$ is $2^{\Theta(\sigma^2)}$, so most rules require $\Omega(\sigma^2)$ bits. Within-level variation in $K(\rho)$ is $O(\sigma \log \sigma)$, sub-leading compared to the $\sigma^2$ inter-level gap. This justifies stating the PRIMO conjecture at the level of $\sigma$ rather than individual $K(\rho)$. See~\cite{kowalczyk2026hierarchy} for the full result.
\end{remark}

\subsection{Open problems}

\begin{enumerate}
\item Construct additional ($\Ineg$, $\Phipos$) witnesses beyond fixed-grid-noise. Conjecture~\ref{conj:conditional} shows that such witnesses must violate at least one of (M1)--(M3'). Candidates include cellular automata on fixed lattices and Dehn-twist constructions on toroidal grids.

\item Extend the computational study to systematic enumeration: all connected rules at signatures $3_2 \to 4_2$ and $4_2 \to 5_2$, random DPO rules at higher signatures, and rules from the Game of Intelligence anti-loop catalog.

\item Characterize the class of $\Ipos$ rules: is it strictly larger than the class of Bayesian graph dynamical systems?

\item Investigate whether the observed frequency asymmetry ($\Ipos$ more common than $\Phipos$) persists under systematic enumeration.
\end{enumerate}


%% ═══════════════════════════════════════════════════════
\section{Conditional implication: $\Phi$-positivity implies I-positivity for growth rules}
\label{sec:conditional}

The $2 \times 2$ classification table (\cref{tab:cells}) shows that the ($\Ineg$, $\Phipos$) cell contains only a single rule (fixed-grid-noise) among the 33-rule test suite. In a separate systematic enumeration of all 16 connected DPO growth rules at signatures $\sigma = 1$ through $\sigma = 4$, the ($\Ineg$, $\Phipos$) cell is entirely empty: every $\Phipos$ DPO rule is also $\Ipos$. This pattern is not a consequence of the predicate definitions---\cref{thm:independence} establishes their logical independence---but reflects a structural property of DPO growth dynamics.

\textbf{Conditions.} A graph dynamical system $(G_0, R, T)$ is a \emph{monotone growth system} if it satisfies:

\begin{itemize}
\item[\textbf{(M1)}] \emph{Monotone growth:} $|V(G_t)|$ and $|E(G_t)|$ are non-decreasing in $t$.
\item[\textbf{(M2)}] \emph{Bounded local modification:} $G_{t+1}$ is obtained from $G_t$ by adding vertices and edges (and possibly rewiring a bounded number of existing edges per matched subgraph), with the total Frobenius-norm perturbation to the adjacency matrix at step $t$ bounded by $\|{\Delta A_t}\|_F \leq C \cdot |E_t|^{1/2}$ for a constant $C$ independent of $t$.
\item[\textbf{(M3')}] \emph{Eventual eigenspace gap stability:} there exists $t_1$ such that for all $t \geq t_1$, the graph $G_t$ has at least $k + 1$ non-zero Laplacian eigenvalues, and the cluster gap---the distance from the $k$-th eigenvalue cluster to the nearest distinct cluster---satisfies $\gap^{\mathrm{cluster}}(L_t) \geq \gamma > 0$.
\end{itemize}

\begin{remark}[On M3']
The cluster gap rather than the raw eigenvalue gap is the correct stability measure. Tree-like growth rules produce graphs with high eigenvalue multiplicity at $\lambda = 1$, giving a raw gap of zero but a large and stable cluster gap. The Davis--Kahan theorem~\cite{davis1970rotation,wedin1972perturbation} applies to eigenspaces, so cluster gap stability suffices.
\end{remark}

\begin{conjecture}[Conditional implication]\label{conj:conditional}
Let $(G_0, R, T)$ be a graph dynamical system satisfying \textup{(M1)}, \textup{(M2)}, and \textup{(M3')}. If $(G_0, R, T)$ is $\Phipos$, then it is $\Ipos$.
\end{conjecture}

The proof has two parts, corresponding to the two gates of the I-predicate.

\subsection*{Compression gate}

\begin{proposition}[Compression gate for deterministic polynomial-growth trajectories]\label{prop:compression}
Let $(G_0, R, T)$ be a DPO graph dynamical system under canonical-ordering GPI satisfying \textup{(M1)}. Suppose the trajectory has polynomial edge growth: there exist constants $c > 0$ and $\alpha > 0$ such that $|E(G_t)| \geq c \cdot t^{\alpha}$ for all $t \geq 1$. Then the trajectory compression ratio satisfies
\[
\compratio = O\!\left(\frac{\log T}{T^{\alpha+1}}\right).
\]
In particular, for any threshold $\rhostar > 0$, there exists $T_0$ depending on $G_0$, $R$, $c$, $\alpha$ such that $\compratio < \rhostar$ for all $T > T_0$.
\end{proposition}

\begin{corollary}
If $(G_0, R, T)$ satisfies \textup{(M1)} and is $\Phipos$, then for $T$ sufficiently large, the compression gate ($\rho < \rhostar = 0.85$) is satisfied.
\end{corollary}

\begin{proof}[Proof of Proposition~\ref{prop:compression}]
\emph{Step~1 (Trajectory determinism).} Under canonical-ordering GPI, the matching at each step is determined by the graph: the set of all injective homomorphisms $L \to G_t$ is enumerable, the maximal independent subset is selected by greedy lexicographic ordering on match vertices, and DPO rewriting of each match is unique (the pushout exists and is unique in $\Graphinj$ for injective rules). Therefore the map $f_R: G_t \mapsto G_{t+1}$ is deterministic, and the entire trajectory $\tau = (G_0, G_1, \ldots, G_T)$ is determined by the triple $(G_0, R, T)$.

\emph{Step~2 (Upper bound on compressed size).} By Step~1, the trajectory is a deterministic function of $(G_0, R, T)$. A compressor that stores the triple $(G_0, R, T)$ plus a fixed decompression program achieves compressed size $|C(S(\tau))| \leq |G_0|_{\text{bits}} + |R|_{\text{bits}} + 2 \log_2 T + c_{\text{sim}}$, where $c_{\text{sim}}$ is a constant independent of $T$. Hence $|C(S(\tau))| = O(\log T)$.

\emph{Step~3 (Lower bound on raw size).} The raw serialized trajectory satisfies $|S(\tau)| = \sum_{t=0}^{T} |S(G_t)| \geq \sum_{t=1}^{T} c \cdot t^{\alpha} = \Omega(T^{\alpha+1})$.

\emph{Step~4 (The ratio).} $\compratio = |C(S(\tau))| / |S(\tau)| \leq O(\log T) / \Omega(T^{\alpha+1}) = O(\log T / T^{\alpha+1})$, which tends to~0 as $T \to \infty$.
\end{proof}

\begin{proof}[Proof of Corollary]
$\Phipos$ requires at least one aggregate quantity to follow a polynomial law with residual below $\lawstar$. Under (M1), this implies polynomial edge growth with $\alpha \geq 1$. The Proposition then applies.
\end{proof}

\subsection*{Embedding convergence}

$\Phipos$ requires stable spectral dimension, which constrains the bulk Laplacian eigenvalue distribution to converge. Under (M1) and~(M2), the per-step embedding perturbation satisfies $\|X_{t+1}^e - X_t^e\|_F \leq C_1 \cdot \|\Delta A_t\|_F$ via the embedding continuity condition. Under (M3'), the Davis--Kahan/Wedin $\sin\theta$ theorem~\cite{davis1970rotation,wedin1972perturbation} gives:
\[
\sin\thetamax\bigl(\col(X_t^e),\; \col(X_{t+1}^e)\bigr) \leq \frac{\|X_{t+1}^e - X_t^e\|_F}{\gap^{\mathrm{cluster}}(L_t)} \leq \frac{C_1 \cdot \|\Delta A_t\|_F}{\gamma}.
\]

As the graph grows under (M1), the relative perturbation $\|\Delta A_t\|_F / \|A_t\|_F$ tends to zero (each step adds a bounded modification to an increasingly large graph). The angular perturbation per step therefore decreases, and the cosine-to-final sequence $c_t^e$ becomes monotonically increasing for $t$ sufficiently large, giving $\taufinal > \taustar = 0.5$.

\begin{remark}
The fixed-grid-noise rule (the sole ($\Ineg$, $\Phipos$) witness from \cref{thm:independence}) violates (M1): the graph has fixed size with one random edge swap per step. The Dehn-twist construction---a toroidal grid with periodic relabeling---provides a sharper theoretical counterexample: it satisfies $\Phipos$ with perfect spectral dimension $\sprdim = 2$, but violates (M1) and produces oscillating embeddings.
\end{remark}

\textbf{Computational verification.} We measured the Laplacian eigenvalue gaps and adjacency singular value gaps for all 16~enumerated DPO growth rules at signatures $\sigma = 1$ through $\sigma = 4$, across all four canonical initial graphs, at $T = 30$ (with $T = 60$ for three borderline rules). Results:

\begin{itemize}
\item All 16 rules satisfy (M1) by construction.
\item Zero eigenvalue crossings were observed across all 16~rules, all 4~seeds, and all time steps.
\item All 9~$\Phipos$ rules satisfy (M3'). Three rules produce tree-like graphs with high eigenvalue multiplicity at $\lambda = 1$; the raw gap is zero but the cluster gap is large (minimum~$0.56$) and stable.
\item Two $\Phineg$ rules also have stable eigenspace gaps but fail $\Phipos$ on spectral dimension instability ($\sigma_{\sprdim} > 0.08$), confirming that (M3') is not the binding constraint.
\end{itemize}

The conjecture explains the empty ($\Ineg$, $\Phipos$) cell in the DPO enumeration as a consequence of monotone growth structure, rather than a failure of predicate independence.


%% ═══════════════════════════════════════════════════════
\appendix

\section{Computational reproducibility}
\label{app:reproducibility}

All computations performed in Python~3.12 with NetworkX, NumPy, SciPy. Random seed fixed at~42. Complete source code: \texttt{primo\_diagnostic\_v5.py} (available as supplementary material). 33~rules, 4~seeds, 3~embeddings, $T = 30$ steps per trajectory. Total runtime: $< 5$~minutes on a single CPU core.

\section{Embedding details}
\label{app:embeddings}

\textbf{Laplacian eigenvectors ($e_L$).} For a graph $G$ with $n$ nodes, compute the graph Laplacian $L = D - A$. The embedding is the matrix of the $d = 5$ leading non-trivial eigenvectors (those corresponding to the smallest $d$ non-zero eigenvalues). For graphs with fewer than $d + 1$ non-zero eigenvalues, the embedding dimension is reduced accordingly.

\textbf{Random projection ($e_R$).} Fix a random Gaussian matrix $M \in \mathbb{R}^{n_{\max} \times d}$ drawn once from $\mathcal{N}(0, 1)$ with random seed~0. For a graph $G$ with $n$ nodes, the embedding is $X = A \cdot M_{n \times d}$ where $A$ is the adjacency matrix and $M_{n \times d}$ is the first $n$ rows and $d = 5$ columns of $M$. The matrix $M$ is fixed across all time steps and all rules.

\textbf{Degree profile ($e_D$).} For a graph $G$ with $n$ nodes, compute for each node $v$ the feature vector: (1)~degree/$\max$ degree, (2)~clustering coefficient, (3)~mean neighbor degree/$\max$ degree, (4)~2-hop neighborhood size/$n$, (5)~shortest-path distance to the highest-degree node/$n$. The embedding is the $n \times 5$ matrix of these features. This embedding has no sign ambiguity and is Lipschitz in the edit distance.

\section{Rule descriptions}
\label{app:rules}

\textbf{Hand-crafted rules.} Identity (do nothing); add random edge (uniform random non-edge); preferential attachment (new node to highest-degree); subdivision (random edge replaced by path of length~2); triangle closure (close first open triangle found); grid growth (rebuild as next-size grid); line growth (extend path by one); progressive compression (merge two lowest-degree nodes); star growth (new leaf to hub~0); cycle-then-fill (build cycle of~10, then fill diagonals); ER random (independent $G(n, 0.3)$ each step); copy-with-noise (copy graph, 20\% chance to remove/add one edge); lattice rewire (grow $4 \times 4$ grid, then degree-preserving random edge swaps); fixed-grid-noise (rebuild $5 \times 5$ grid each step, swap one edge); sorting edges (rewire highest-$|\Delta\text{deg}|$ edge to reduce degree variance).

\textbf{Catalog rules.} From the complete catalog of GPI rewrite rules: vertex sprouting (every vertex sprouts a leaf); one-sided edge sprouting (matched edge kept, fresh vertex to one endpoint); triangle completion (matched edge plus fresh vertex to both endpoints); edge deletion (remove matched edge); edge rewiring (delete matched edge, create fresh vertex connected to one endpoint).

\textbf{Standard graph families.} Barab\'asi--Albert ($m = 2$ new edges per step); Watts--Strogatz (grow ring lattice, rewiring probability~$0.1$); caterpillar (extend spine, sprout leaf); complete bipartite (add to smaller partition, connect to all of larger); degree regularization (add node to lowest-degree, rewire if any node exceeds $2\times$ mean degree).

\textbf{Random DPO rules.} Five rules generated by: LHS $= \Kn{2}$, RHS $= 3$ vertices with random edge subset (each of 3 possible edges included with probability~$0.5$), interface map $=$ random injection from $\{0,1\}$ to $\{0,1,2\}$. Seeds 100--104.

\textbf{Candidate witnesses.} Hierarchical tree (split first leaf into two children); hub-sort (rewire lowest-degree node toward hub's neighborhood); encode-compress (move edge from highest-above-mean to lowest-below-mean degree node).


%% ═══════════════════════════════════════════════════════
\begin{acknowledgments}
This work is part of the PRIMO research programme. Code and data are available at \url{https://github.com/KarolFilipKowalczyk/PRIMO}.
\end{acknowledgments}

\bibliographystyle{plain}
\bibliography{../shared/primo}

\end{document}
